% Front matter and introduction. Compile the repository root document instead.
% !TeX root = Stratified_Relational_Workspace_Theory.tex
\pagenumbering{gobble}

\begin{titlepage}
\thispagestyle{empty}
\centering
\vspace*{1.2in}
{\Huge\bfseries\color{PartColor} Stratified Relational\\[0.18em] Workspace Theory\par}
\vspace{0.55in}
{\Large A Core--Transfer Framework for Distributed Cognition,\\ Collaboration, and Epistemic Systems\par}
\vspace{0.8in}
{\large Komorebi P.\par}
\vspace{0.65in}
{\large\bfseries Foundations, Dynamics, Agency, and Empirical Certification\par}
\vfill
{\color{RuleColor}\rule{0.70\textwidth}{0.8pt}\par}
\vspace{0.25in}
{\small August 2026\par}
\end{titlepage}
\clearpage
\pagenumbering{roman}

\begin{abstract}
Collaborative reasoning increasingly spans people, computational models, documents, software,
instruments, and persistent records. Yet the existence of many interacting components does not,
by itself, establish a shared workspace, and successful group performance does not show that the
operative state is genuinely distributed rather than copied, centralized, or merely correlated.
Stratified Relational Workspace Theory develops a formal answer to this problem. A realization
is represented by heterogeneous local state spaces connected by transfer maps. Its exact shared
core is the space of global compatible sections, equivalently the fixed subspace of relational
holonomy when the transfers are invertible. A transfer Laplacian identifies approximate shared
directions and provides spectral stability, gauge covariance, and pathwise leakage bounds.

The structural theory is joined to a dynamic and causal account. Canonical predictive states
determine whether a proposed workspace retains the conditional law relevant to a declared
future task. Feasible interventions, coalition risk, matched relational nulls, and reconstruction
tests distinguish causally operative distributed organization from archives, replicated memory,
and centralized controllers. Further results address persistence under changing realizations,
stable reconstruction after component loss, identity under replacement, local-to-global gluing,
hierarchical composition, learned transfer structure, boundary recovery, goal-conditioned action,
and certificate-directed repair. Because distributedness can be destroyed by merging components
or manufactured by splitting a centralized state into complementary shares, the theory also treats
decomposition as part of the claim. Faithful coarse-graining preserves exact structural cores, while
family-robust certification requires every gate to retain a positive margin across a declared set of
admissible decompositions. The same discipline is applied to the complete scientific specification.
Targets, losses, horizons, interventions, nulls, and thresholds are made explicit, and the theory develops
specification envelopes, transport radii, experimental-equivalence classes, partial-identification bounds,
and uniform finite-sample certification. A certificate is therefore distinguished from a favorable
analysis: it may be pointwise valid, specification-robust, sign-identified, or statistically unresolved.
The resulting rule is noncompensatory. Structural coherence, predictive alignment, causal efficacy,
relational irreducibility, non-collapse, reconstructive continuity, provenance, and self-maintenance are
separate claims, and failure of one cannot be hidden by success on another.

Exact finite examples show how the same framework applies to human collaboration, human--AI
research loops, and multi-agent control. A review of adjacent work locates the individual
mathematical ingredients in established traditions including distributed cognition, local systems
and sheaves, spectral perturbation, sufficient statistics, structural causal models, frame reconstruction,
system identification, and feedback control. The candidate contribution is therefore
not a new name for any one of these constructions. It is their integration into a single falsifiable
account of persistent relational workspaces and the matched near misses that such an account
must reject. The framework does not infer phenomenal consciousness, personhood, or moral
status from coherence, agency, or self-maintenance.
\end{abstract}

\tableofcontents
\clearpage
\pagenumbering{arabic}

\section*{Introduction}
\addcontentsline{toc}{section}{Introduction}

Scientific and technical work rarely resides in a single mind or device. A proof may be divided
among a conjecture, a critic's counterexample, a notebook, a computer algebra system, and a
verification script. A human--AI research process may distribute objectives, formal manipulations,
execution, and persistent memory across different components. Multi-agent systems divide sensing,
planning, communication, and repair in still other ways. In each case, it is natural to speak of a
shared workspace. The phrase is useful, but without a formal criterion it can conceal several
fundamentally different organizations.

A common archive is not necessarily a workspace. Nor is a group whose apparent coordination is
produced by one complete controller surrounded by decorative channels. Two systems can contain
the same local information while differing in how that information is related, transported, and used.
Likewise, a system can perform well without preserving its operative organization through interruption,
replacement, or component loss. The central problem is therefore not simply whether information is
shared. It is whether there is a persistent relational state whose structure and causal capacities cannot
be reduced to any single declared component.

The philosophical background is supplied by extended and distributed cognition, transactive memory,
collective intelligence, and theories of group agency \cite{Hutchins1995,ClarkChalmers1998,Wegner1987,ListPettit2011}.
Those traditions motivate the possibility that cognitive organization may span people and artifacts.
The theory asks a narrower mathematical question: what would have to be measured and proved
before such an organization should be called a relational workspace? The answer must distinguish
physical realization from abstract identity, local representation from globally compatible content,
and statistical association from causal operation.

The static theory begins with an interaction graph whose vertices carry heterogeneous local state spaces
and whose edges carry transfer maps. Each local state is decomposed into a transferable core and a
substrate-specific residual. Exact shared content is represented by global compatible sections. When
core transfers are invertible, the same object appears as the fixed subspace of the holonomy representation.
Approximate shared content is extracted as a low-energy spectral subspace of a transfer Laplacian.
This formulation supplies explicit obstruction, stability, gauge, and leakage calculations rather than
relying on an informal appeal to shared context.

Structure alone is insufficient. A proposed workspace must also retain the predictive state relevant to
a declared future, respond under feasible intervention, outperform matched relational nulls, and resist
reduction to an individual node or a local-storage surrogate. Persistence introduces further requirements:
the workspace must remain trackable as its realization changes, preserve nontrivial content, support
reconstruction after loss, and maintain an identity profile under admissible replacement. Local workspaces
must glue without hidden incompatibility, and hierarchical organization must not conceal centralization.
Formation claims are evaluated separately from operational claims so that a scaffolded but functional
workspace is not dismissed, while an endogenously learned but causally irrelevant pattern is not promoted.
A separate decomposition analysis asks which claims survive faithful coarse-graining and plausible
alternative interface descriptions. This is necessary because arbitrary merging can conceal relationality,
while arbitrary splitting can manufacture it. The analysis is then extended from decomposition to the full
scientific specification. A workspace claim depends on its declared target, loss, horizon, intervention
family, null panel, and thresholds. The theory therefore distinguishes pointwise validity from robustness
across defensible specifications and distinguishes both from identifiability under the experiments that are
actually available. This prevents a favorable target or threshold from being mistaken for a general result
and prevents an observationally indistinguishable hidden controller from being ruled out by assertion.

The theory extends to domain-relative agency only after these lower-level conditions are in place. Goals
must persist without repeated external rewriting, counterfactual values must guide feasible actions, damage
must be detected through the workspace, and repairs must improve a matched no-repair trajectory while
preserving identity. These criteria define an engineering notion of workspace agency. They neither assert
nor deny phenomenal consciousness. Consciousness, personhood, and moral status remain separate questions
that require independent evidence.

Most of the mathematical tools used here are established. Inverse limits, holonomy, connection Laplacians,
perturbation theory, sufficient statistics, causal models, frame reconstruction, system identification, and
feedback control all have substantial prior literatures. The contribution sought here is integrative: a
noncompensatory framework that places these tools in one common architecture and specifies the near-miss
systems that each layer must reject. Exact examples provide calibration, while the motivating human--AI
collaboration is treated conservatively as a candidate case rather than as a certified instance.

The paper proceeds from exact structure to approximate structure, dynamic operation, relational
irreducibility, persistence, reconstruction, identity, gluing, formation, boundary discovery, agency,
decomposition robustness, specification robustness, identifiability, and empirical certification. It then
gives worked cases, situates the framework relative to adjacent work, and closes with its philosophical
consequences, limitations, and open problems.

\section*{Notation and Conventions}
\addcontentsline{toc}{section}{Notation and Conventions}

A concrete realization is denoted by \(\mathfrak R\). Its interaction graph is \(G=(V,E,s,t)\), with local
state spaces \(X_i\) at vertices and transfer maps \(T_e:X_{s(e)}\to X_{t(e)}\) on directed edges. When a
core--residual decomposition is used, \(X_i=C_i\oplus R_i\), and \(H_e\) denotes the induced core transfer.
The exact structural workspace is the compatible-section space \(\Gamma(H)\); the transfer discrepancy is
\(D_H\), the transfer Laplacian is \(L_H=D_H^*D_H\), and \(\Gamma_\varepsilon(H)\) denotes a selected
low-energy workspace subspace. Dynamic workspace states are written \(W_t\), while \(Z_t\) denotes the
canonical predictive state associated with a declared future variable \(Y_t^+\).

Unless stated otherwise, the structural results use finite-dimensional real or complex Hilbert spaces and
finite graphs. Probabilistic statements are made on standard Borel spaces when regular conditional laws are
required. All approximate conclusions are indexed by explicit tolerances, spectral gaps, intervention
families, loss functions, or null classes. A claim of relational irreducibility is always relative to a declared
decomposition of the realization into components. The symbol \(\mathcal D\) denotes one such
decomposition, and \(\mathfrak D\) denotes an admissible family used for robustness analysis. The complete
scientific specification is
\[
\varsigma=(\mathcal D,Y^+,\ell,\tau,\mathcal I,\mathcal N,\Theta),
\]
and \(\mathfrak S\) denotes an admissible family of specifications. A gate margin is written
\(m_q(W;\varsigma)\), with positive sign indicating passage under the declared direction and threshold.

The terms \emph{workspace}, \emph{agency}, \emph{consciousness}, \emph{personhood}, and \emph{moral status}
are not interchangeable. A structural workspace may fail every operational test. An operational workspace
may lack persistence. A persistent relational workspace may lack goals or repair. A workspace agent, as
defined here, is a domain-relative functional organization and is not thereby a conscious subject. Definitions,
propositions, theorems, conjectures, and empirical tests are identified by their ordinary mathematical form;
all proved statements remain conditional on the assumptions given with them.
